\documentclass{amsart}
\usepackage{amsmath,amssymb,amsthm,hyperref,mathtools}

\newtheorem{theorem}{Theorem}
\newtheorem{lemma}{Lemma}
\newtheorem{proposition}{Proposition}
\newtheorem{corollary}{Corollary}
\theoremstyle{definition}
\newtheorem{definition}{Definition}
\newtheorem{remark}{Remark}
\newtheorem{example}{Example}

\DeclareMathOperator{\uce}{uce}
\DeclareMathOperator{\HC}{HC}
\DeclareMathOperator{\im}{im}
\DeclareMathOperator{\Hom}{Hom}
\DeclareMathOperator{\id}{id}
\DeclareMathOperator{\Aut}{Aut}
\newcommand{\g}{\mathfrak{g}}
\newcommand{\h}{\mathfrak{h}}
\newcommand{\nn}{\mathfrak{n}}
\newcommand{\sw}{\mathfrak{sl}}
\newcommand{\Om}{\Omega}
\newcommand{\Gm}{\Gamma}
\newcommand{\ra}{\rightarrow}

\begin{document}

\title{Universal central extensions of equivariant map algebras}
\maketitle

Throughout, $k$ is a field of characteristic $0$, $\g$ is a finite-dimensional
Lie algebra over $k$, $A$ is a unital commutative associative $k$-algebra, and
$\Gm$ is a finite group acting by Lie algebra automorphisms on $\g$ and by algebra
automorphisms on $A$. We let $\Gm$ act diagonally on the current algebra
$L:=\g\ot A$, and set
\[
   M=L^{\Gm}=(\g\ot A)^{\Gm}=\{u\in\g\ot A:\gamma\cdot u=u\ \text{for all }\gamma\in\Gm\}.
\]
We assume $M$ is perfect.  All tensor products, $\Hom$'s, exterior and symmetric
powers are over $k$ unless decorated otherwise.

Because $\operatorname{char}k=0$ and $\Gm$ is finite, the averaging idempotent
\[
   e=\frac{1}{|\Gm|}\sum_{\gamma\in\Gm}\gamma\in k[\Gm]
\]
acts on every $k[\Gm]$-module $V$ as a projector with image $V^{\Gm}$, giving a
canonical $\Gm$-stable decomposition $V=V^{\Gm}\oplus(1-e)V$.  Hence the functor
$(-)^{\Gm}$ is \emph{exact} on $k[\Gm]$-modules (it is the direct summand $e$ of
the identity functor), and $e$ commutes with every $\Gm$-equivariant linear map.
We use these two facts constantly.

\paragraph{Summary of results.}  Part (a) is settled completely and unconditionally
(Theorem~\ref{thm:partA}): $\uce(M)=\Lambda^2M/J_M$ with centre $H_2(M)$.  For part
(b) we prove (Proposition~\ref{prop:snake}) that the natural comparison map $\Phi$
is an isomorphism \emph{iff} an explicit linear map
$\phi\colon H_2(M)\to H_2(\g\ot A)^{\Gm}$ is.  We then determine $\phi$.  When the
isotypic complement $\nn=(1-e)L$ is a Lie ideal we compute $\phi$ exactly: it is
always injective, with $\operatorname{coker}\phi\cong H_2(\nn)^{\Gm}$, hence an
isomorphism \emph{iff} $H_2(\nn)^{\Gm}=0$ (Theorem~\ref{thm:partB}(2)).  For $\g$
semisimple, $\phi$ is an isomorphism for every finite $\Gm$ and every $A$ with $M$
perfect (Theorem~\ref{thm:partB}(4); this is the theorem of Neher--Savage, whose
hypotheses we verify and whose mechanism we describe through an explicit
$\Gm$-invariant Killing contraction, Lemma~\ref{lem:kappa}).  Finally,
$\phi$ can fail to be surjective when $\g$ is merely perfect with $H_2(\g)\ne0$;
we give a completely explicit verified counterexample (Theorem~\ref{thm:partB}(5),
Example~\ref{ex:fail}).

\bigskip
\section{Universal central extensions: the machinery}

\begin{definition}\label{def:uce}
For a Lie algebra $L$ over $k$, a \emph{central extension} is a surjective Lie
homomorphism $\pi\colon \widehat L\ra L$ with $\ker\pi\subseteq Z(\widehat L)$. It
is \emph{universal} if for every central extension $\pi'\colon E\ra L$ there is a
unique Lie homomorphism $\varphi\colon\widehat L\ra E$ with
$\pi'\circ\varphi=\pi$.
\end{definition}

\begin{theorem}[Garland--Kassel]\label{thm:uce-exists}
Let $L$ be a perfect Lie algebra over $k$.  Put
\[
   \uce(L)=\Lambda^2L/J,\quad
   J=\operatorname{span}_k\{[x,y]\wedge z+[y,z]\wedge x+[z,x]\wedge y:x,y,z\in L\},
\]
write $\langle x,y\rangle$ for the image of $x\wedge y$, define the bracket
$[\langle x,y\rangle,\langle u,v\rangle]=\langle[x,y],[u,v]\rangle$ and the map
$\pi\colon\uce(L)\ra L$, $\langle x,y\rangle\mapsto[x,y]$.  Then $\pi$ is the
universal central extension of $L$, the algebra $\uce(L)$ is perfect, and
\[
   \ker\pi=Z(\uce(L))\cong H_2(L):=H_2(L;k),
\]
the second Chevalley--Eilenberg homology with trivial coefficients.  Concretely,
with the boundary maps
\[
   \Lambda^3L\xrightarrow{d_3}\Lambda^2L\xrightarrow{d_2}L,\quad
   d_2(x\wedge y)=[x,y],\
   d_3(x\wedge y\wedge z)=[x,y]\wedge z-[x,z]\wedge y+[y,z]\wedge x,
\]
one has $H_2(L)=\ker d_2/\im d_3=\ker d_2/J$, equal to $\ker\pi\subseteq\Lambda^2L/J$.
\end{theorem}

\begin{proof}
Classical (Weibel, \emph{Introduction to Homological Algebra}, \S7.9;
Kassel--Loday, \emph{Ann. Inst. Fourier} 32 (1982); Garland, \emph{Publ. IH\'ES} 52
(1980)).  We record the points used.

\emph{Bracket well defined.}  The bilinear map
$(x\wedge y,u\wedge v)\mapsto[x,y]\wedge[u,v]$ on $\Lambda^2L$ kills $J$ in each
slot: for $j=[x,y]\wedge z+[y,z]\wedge x+[z,x]\wedge y$,
\[
   [j,u\wedge v]\mapsto\bigl([[x,y],z]+[[y,z],x]+[[z,x],y]\bigr)\wedge[u,v]=0
\]
by the Jacobi identity, and symmetrically; Jacobi for $\uce(L)$ reduces to Jacobi
for $L$.

\emph{$\pi$ central, $\uce(L)$ perfect.}  $\pi$ is a surjective homomorphism since
$L=[L,L]$; if $[x,y]=0$ then $[\langle x,y\rangle,\langle u,v\rangle]=
\langle0,[u,v]\rangle=0$, so $\ker\pi$ is central.  A generator $\langle
x,y\rangle$ with $y=\sum_i[c_i,d_i]$ equals $\sum_i[\langle x,c_i\rangle,\langle
e_i,f_i\rangle]$ ($[e_i,f_i]=d_i$), so $\uce(L)=[\uce(L),\uce(L)]$.

\emph{Universality.}  For a central extension $\pi'\colon E\ra L$ with $k$-linear
section $s$, set $\varphi\langle x,y\rangle=[s(x),s(y)]$; centrality of $\ker\pi'$
makes this independent of $s$ and kills $J$, giving a Lie homomorphism with
$\pi'\varphi=\pi$.  Two lifts of $\pi$ differ by a central map vanishing on
$\uce(L)=[\uce(L),\uce(L)]$, hence coincide.

\emph{Homology.}  For a perfect central extension, $Z(\widehat L)=\ker\pi$, and
$\ker\pi=\ker d_2/\im d_3=H_2(L)$ by the displayed Chevalley--Eilenberg complex.
\end{proof}

\begin{corollary}[Functoriality and induced group actions]\label{cor:funct}
$\uce$ is a functor on perfect Lie algebras: $f\colon L\ra L'$ induces the unique
lift $\uce(f)\colon\uce(L)\ra\uce(L')$,
$\langle x,y\rangle\mapsto\langle f(x),f(y)\rangle$, with
$\pi'\circ\uce(f)=f\circ\pi$.  A $\Gm$-action on $L$ by automorphisms therefore
lifts to the action $\gamma\cdot\langle x,y\rangle=\langle\gamma x,\gamma
y\rangle$ on $\uce(L)$, and $\pi$ is $\Gm$-equivariant; this is the unique lift
through $\pi$, settling the lifting assertion of part (b).
\end{corollary}
\begin{proof}
$\langle x,y\rangle\mapsto\langle f(x),f(y)\rangle$ kills $J$ and respects
brackets; uniqueness follows as in Theorem~\ref{thm:uce-exists}.  Functoriality is
immediate from the generator formula, so a group action lifts to a group action.
Equivariance: $\pi(\gamma\langle x,y\rangle)=[\gamma x,\gamma y]=\gamma[x,y]$.
\end{proof}

\section{The comparison map and its homological control}

\begin{lemma}[The two central extensions of $M$]\label{lem:two-ext}
The $\Gm$-equivariant projection $\pi\colon\uce(L)\ra L$ restricts to a central
extension
$\pi^{\Gm}\colon \uce(L)^{\Gm}\ra M=L^{\Gm}$, with
$\ker\pi^{\Gm}=(\ker\pi)^{\Gm}=H_2(L)^{\Gm}$ central in $\uce(L)^{\Gm}$.  As $M$ is
perfect, the universal property of $\uce(M)$ supplies a unique Lie homomorphism
\[
   \Phi\colon\uce(M)\longrightarrow\uce(L)^{\Gm},\qquad \pi^{\Gm}\circ\Phi=\pi_M,
\]
given on generators by $\Phi\langle x,y\rangle=\langle x,y\rangle$ ($x,y\in M$).
This is the natural comparison map of the problem.
\end{lemma}
\begin{proof}
$\Gm$-equivariance (Corollary~\ref{cor:funct}) and exactness of $(-)^{\Gm}$ give
surjectivity of $\pi^{\Gm}$: for $w\in M$ pick $\widetilde w$ with
$\pi\widetilde w=w$; then $e\widetilde w\in\uce(L)^{\Gm}$ and
$\pi(e\widetilde w)=ew=w$.  $\ker\pi^{\Gm}=\ker\pi\cap\uce(L)^{\Gm}=
(\ker\pi)^{\Gm}=H_2(L)^{\Gm}$ by exactness of $e$, central because
$\ker\pi=Z(\uce(L))$.  Thus $\pi^{\Gm}$ is a central extension of $M$; $\Phi$
exists and is unique by Theorem~\ref{thm:uce-exists}, and the generator formula
defines a lift of $\pi_M$, hence equals $\Phi$.
\end{proof}

\begin{lemma}[The invariant subcomplex computes invariant homology]
\label{lem:invariant-complex}
For any $\Gm$-Lie algebra $L$ ($\Gm$ finite, $\operatorname{char}k=0$), the
Chevalley--Eilenberg complex $C_\bullet(L)=\Lambda^\bullet L$ is a complex of
$k[\Gm]$-modules, $e$ is a chain map, and
$H_n\bigl(C_\bullet(L)^{\Gm}\bigr)\cong H_n(L)^{\Gm}$ for all $n$.
\end{lemma}
\begin{proof}
$\Gm$ acts functorially on $\Lambda^\bullet L$ and commutes with $d$ (built from
the $\Gm$-equivariant bracket), so $e$ is a chain map and $C_\bullet(L)=
C_\bullet(L)^{\Gm}\oplus(1-e)C_\bullet(L)$ as complexes.  Since $e$ is an exact
projector commuting with $d$,
\[
   H_n(C_\bullet(L)^{\Gm})=\frac{e\ker d}{e\im d}=e\,\frac{\ker d}{\im d}
   =H_n(L)^{\Gm}.\qedhere
\]
\end{proof}

\begin{proposition}[$\Phi$ is controlled by an explicit homology map]
\label{prop:snake}
Let $\phi:=\Phi|_{H_2(M)}\colon H_2(M)\ra H_2(L)^{\Gm}$ (it lands in
$\ker\pi^{\Gm}=H_2(L)^{\Gm}$ since $\pi^{\Gm}\Phi=\pi_M$).  Then:
\begin{enumerate}
\item[(a)] $\ker\Phi\cong\ker\phi$ and $\operatorname{coker}\Phi\cong
\operatorname{coker}\phi$; in particular $\Phi$ is injective (resp.\ surjective,
resp.\ iso) iff $\phi$ is.
\item[(b)] $\phi$ is the map on $H_2$ induced by the chain inclusion
\[
   C_\bullet(M)=\Lambda^\bullet(L^{\Gm})\ \hookrightarrow\ (\Lambda^\bullet L)^{\Gm}
   =C_\bullet(L)^{\Gm},
\]
followed by the identification $H_2(C_\bullet(L)^{\Gm})=H_2(L)^{\Gm}$ of
Lemma~\ref{lem:invariant-complex}.
\end{enumerate}
\end{proposition}
\begin{proof}
(a) By Lemma~\ref{lem:two-ext} there is a commutative diagram of short exact
sequences with identity on the right:
\[
\begin{array}{ccccccccc}
0 &\!\!\ra\!\!& H_2(M) &\!\!\ra\!\!& \uce(M) &\!\!\xrightarrow{\pi_M}\!\!& M &\!\!\ra\!\!& 0\\[2pt]
  &   & \big\downarrow{\scriptstyle\phi} & & \big\downarrow{\scriptstyle\Phi} & & \big\|\ & & \\[2pt]
0 &\!\!\ra\!\!& H_2(L)^{\Gm} &\!\!\ra\!\!& \uce(L)^{\Gm} &\!\!\xrightarrow{\pi^{\Gm}}\!\!& M &\!\!\ra\!\!& 0.
\end{array}
\]
The snake lemma (right map an isomorphism) gives
$0\ra\ker\phi\ra\ker\Phi\ra0\ra\operatorname{coker}\phi\ra
\operatorname{coker}\Phi\ra0$, so $\ker\Phi\cong\ker\phi$ and
$\operatorname{coker}\Phi\cong\operatorname{coker}\phi$.

(b) $\Phi\langle x,y\rangle=\langle x,y\rangle$ ($x,y\in M$) is induced by
$\Lambda^2M\hookrightarrow\Lambda^2L$ on bracket-kernels modulo the $J$'s; on
homology this is exactly the map induced by $\Lambda^\bullet M\hookrightarrow
(\Lambda^\bullet L)^{\Gm}$ composed with Lemma~\ref{lem:invariant-complex} (the
image is $\Gm$-invariant since elements of $M$ are fixed).
\end{proof}

Thus the whole of part (b) is the single explicit map
\begin{equation}\label{eq:phi}
   \phi\colon\ H_2\bigl(\Lambda^\bullet(L^{\Gm})\bigr)\ \longrightarrow\
   H_2\bigl((\Lambda^\bullet L)^{\Gm}\bigr)=H_2(L)^{\Gm},
\end{equation}
induced by $\Lambda^\bullet(L^{\Gm})\hookrightarrow(\Lambda^\bullet L)^{\Gm}$.
The source uses exterior powers of fixed points; the target uses fixed points of
exterior powers.

\section{Part (a): central extensions of $M$ and $\uce(M)$}

\begin{theorem}[Answer to (a); unconditional]\label{thm:partA}
Let $M=(\g\ot A)^{\Gm}$ be perfect.
\begin{enumerate}
\item[(i)] The universal central extension is
\[
   \uce(M)=\Lambda^2M/J_M,\quad
   J_M=\operatorname{span}_k\{[x,y]\wedge z+[y,z]\wedge x+[z,x]\wedge y:x,y,z\in M\},
\]
with $[\langle x,y\rangle,\langle u,v\rangle]=\langle[x,y],[u,v]\rangle$ and
$\pi_M\langle x,y\rangle=[x,y]$; its centre equals its kernel,
\[
   Z(\uce(M))=\ker\pi_M\cong H_2(M)=H_2\bigl((\g\ot A)^{\Gm}\bigr),
\]
the middle homology of $\Lambda^3M\xrightarrow{d_3}\Lambda^2M\xrightarrow{d_2}M$.
\item[(ii)] Up to equivalence, central extensions of $M$ form the vector space
$H^2(M;k)\cong\Hom_k(H_2(M),k)$; each is a unique pushout of the universal one
along a linear map $H_2(M)\ra\mathfrak z$ into its central kernel, the universal
extension corresponding to $\id_{H_2(M)}$.
\end{enumerate}
\end{theorem}
\begin{proof}
(i) is Theorem~\ref{thm:uce-exists} with $L=M$ (legitimate since $M$ is perfect by
hypothesis).  (ii): for a perfect Lie algebra in characteristic $0$, equivalence
classes of central extensions form a torsor over $H^2(M;k)=\Hom(H_2(M),k)$ (the
Chevalley--Eilenberg $H^2$ is the $k$-dual of $H_2$), and the universal extension
is the initial object, obtained as $\id_{H_2(M)}$; every extension is a unique
pushout of it.
\end{proof}

To make $H_2(M)$ \emph{explicit} we combine the structural reduction of part (b)
with Kassel's computation of $H_2(\g\ot A)$.

\begin{theorem}[Kassel]\label{thm:kassel}
Let $\g$ be finite-dimensional over $k$ ($\operatorname{char}k=0$), $A$
commutative unital, $\Om^1_A=\Om^1_{A/k}$ the K\"ahler differentials,
$\HC_1(A)=\Om^1_A/dA$ the first cyclic homology.  If $\g$ is perfect there is a
natural ($\g$- and $\Aut(\g)\!\times\!\Aut(A)$-equivariant) isomorphism
\begin{equation}\label{eq:kassel}
   H_2(\g\ot A)\ \cong\ \bigl(H_2(\g)\ot A\bigr)\ \oplus\
   \bigl(W(\g)\ot\Om^1_A/dA\bigr),
\end{equation}
where $W(\g)=S^2\g/\g\!\cdot\!S^2\g$ is the coinvariant space of the symmetric
square (the universal recipient of $\g$-invariant symmetric bilinear forms).  For
$\g$ simple, $\dim(S^2\g^*)^{\g}=1$ (Killing form), so $W(\g)\cong k$.
\end{theorem}
\begin{proof}
Kassel, \emph{J. Pure Appl. Algebra} 34 (1984) 265--275 (building on
Kassel--Loday).  Structure: $\g\ot\g=\Lambda^2\g\oplus S^2\g$; the $2$-cycles of
$\g\ot A$ split accordingly.  The $\Lambda^2\g$-part contributes $H_2(\g)\ot A$
modulo Jacobi relations; the $S^2\g$-part pairs with the antisymmetric part of
$A\ot A$ and the $d_3$-relations realise it as $\Om^1_A/dA$ (via $a\,db$, with
$a\,db+b\,da=d(ab)\equiv0$), the adjoint action projecting $S^2\g$ onto $W(\g)$.
Naturality in $\g,A$ is part of the statement.
\end{proof}

\begin{corollary}[Simple $\g$]\label{cor:simple}
For $\g$ simple: $H_2(\g)=0$, $W(\g)=k$, and
\[
   H_2(\g\ot A)\cong\Om^1_A/dA,\qquad
   \uce(\g\ot A)=(\g\ot A)\oplus(\Om^1_A/dA),
\]
with $[x\ot a,y\ot b]=[x,y]\ot ab+(x,y)_\kappa\,\overline{a\,db}$,
$(\cdot,\cdot)_\kappa$ the Killing form.  Automorphisms preserve $\kappa$
($\kappa(\gamma x,\gamma y)=\kappa(x,y)$), so $\Gm$ acts trivially on $W(\g)=k$
and only through $\Om^1_A/dA$.  Combined with Theorem~\ref{thm:partB}(4) (which
gives $\uce(M)\cong\uce(\g\ot A)^{\Gm}$ for $\g$ semisimple) and exactness of
$(-)^{\Gm}$:
\[
   \boxed{\ \uce\bigl((\g\ot A)^{\Gm}\bigr)\ \cong\ M\ \oplus\ (\Om^1_A/dA)^{\Gm}\ }
   \qquad(\g\ \text{simple}),
\]
a central extension whose centre is $(\Om^1_A/dA)^{\Gm}$, with $\Gm$ acting by
$\gamma\cdot\overline{a\,db}=\overline{(\gamma a)\,d(\gamma b)}$.
\end{corollary}
\begin{proof}
$H_2(\g)=0$ (Whitehead) and $\dim(S^2\g^*)^\g=1$ (simplicity) give the first line;
$\kappa$-invariance is immediate from $\operatorname{ad}(\gamma x)=\gamma\,
\operatorname{ad}(x)\,\gamma^{-1}$ and the trace formula for $\kappa$.  The boxed
formula combines Theorem~\ref{thm:partB}(4) with
$\uce(\g\ot A)^{\Gm}=(M)\oplus(\Om^1_A/dA)^{\Gm}$, where the splitting of
invariants is the exact functor $(-)^{\Gm}$ applied to
$\uce(\g\ot A)=(\g\ot A)\oplus(\Om^1_A/dA)$.
\end{proof}

\section{Part (b): when is the comparison map an isomorphism?}

By Whitehead's lemmas, $\g$ semisimple has $H_1(\g)=H_2(\g)=0$, so Kassel
\eqref{eq:kassel} becomes the natural isomorphism
\begin{equation}\label{eq:ss-kassel}
   H_2(\g\ot A)\ \cong\ W(\g)\ot\Om^1_A/dA ,\qquad
   W(\g)=\textstyle\bigoplus_{s}k\,\kappa_s ,
\end{equation}
the sum over the simple ideals $\g_s$ of $\g$, with $\g\ot A$ acting trivially and
$\Gm$ acting diagonally (permuting the simple factors $\g_s$ -- equivalently the
$\kappa_s$ -- through its action on $\g$, and through $A$ on $\Om^1_A/dA$).  For
a $\kappa_s$-dual pair of bases $\{p^{(s)}_\alpha\},\{p^{(s)\alpha}\}$ of $\g_s$
(so $\kappa_s(p^{(s)}_\alpha,p^{(s)\beta})=\delta_\alpha^\beta$), the class of
$\overline{\kappa_s\ot a\,db}$ is represented by the $2$-cycle
\begin{equation}\label{eq:cas-cycle}
   z_s(a,b)=\sum_\alpha (p^{(s)}_\alpha\ot a)\wedge(p^{(s)\alpha}\ot b)
   \ \in\ \Lambda^2(\g\ot A),
\end{equation}
which is a cycle because $\sum_\alpha[p^{(s)}_\alpha,p^{(s)\alpha}]=0$ (the
Casimir element of $\g_s$ is central).

\begin{lemma}[The $\Gm$-invariant Killing contraction]\label{lem:kappa}
Let $\g$ be semisimple, $\kappa$ its Killing form (restricting to $\kappa_s$ on the
simple ideal $\g_s$), and $W(\g)=\bigoplus_s k\,\kappa_s$.  Define the linear map
\[
   c\colon\ \Lambda^2(\g\ot A)\longrightarrow W(\g)\ot\Om^1_A/dA,\qquad
   c\bigl((x\ot a)\wedge(y\ot b)\bigr)=\textstyle\sum_s\kappa_s(x,y)\,\kappa_s\ot\overline{a\,db}.
\]
Then:
\begin{enumerate}
\item[(i)] $c$ is $\Gm$-equivariant for the $\Gm$-actions described after
\eqref{eq:ss-kassel};
\item[(ii)] $c\circ d_3=0$, so $c$ descends to $\bar c\colon H_2(\g\ot A)\ra
W(\g)\ot\Om^1_A/dA$;
\item[(iii)] $\bar c$ is an isomorphism, equal to the inverse of the Kassel
isomorphism \eqref{eq:ss-kassel} up to the positive scalars $\dim\g_s$.
\end{enumerate}
\end{lemma}
\begin{proof}
(i) For $\gamma\in\Gm$, $\gamma$ permutes the simple ideals; if
$\gamma\g_s=\g_{s'}$ then $\kappa_{s'}(\gamma x,\gamma y)=\kappa_s(x,y)$ because the
Killing form is $\Aut(\g)$-invariant and a sum over simple ideals.  Hence
\[
   c\bigl(\gamma\cdot((x\ot a)\wedge(y\ot b))\bigr)
   =\sum_s\kappa_s(\gamma x,\gamma y)\,\kappa_s\ot\overline{(\gamma a)\,d(\gamma b)}
   =\sum_{s}\kappa_{s}(x,y)\,(\gamma\kappa_{s})\ot\overline{(\gamma a)\,d(\gamma b)}
   =\gamma\cdot c\bigl((x\ot a)\wedge(y\ot b)\bigr),
\]
where in the middle step we reindexed $s\mapsto s'=\gamma s$ and used
$\gamma\cdot\kappa_s=\kappa_{\gamma s}$.

(ii) Fix simple ideal $\g_s$ and put $T_s(x,y,z)=\kappa_s([x,y],z)$.  Invariance of
$\kappa_s$ ($\kappa_s([u,v],w)=-\kappa_s(v,[u,w])$) makes $T_s$ totally
antisymmetric.  Using $[x\ot a,y\ot b]=[x,y]\ot ab$,
\[
\begin{aligned}
c\,d_3\bigl((x\ot a)\wedge(y\ot b)\wedge(z\ot c)\bigr)
 &=c\Bigl(([x,y]\ot ab)\wedge(z\ot c)-([x,z]\ot ac)\wedge(y\ot b)\\
 &\hspace{2.6em}+([y,z]\ot bc)\wedge(x\ot a)\Bigr)\\
 &=\sum_s\kappa_s\ot\Bigl(\kappa_s([x,y],z)\,\overline{(ab)\,dc}
   -\kappa_s([x,z],y)\,\overline{(ac)\,db}\\
 &\hspace{6.0em}+\kappa_s([y,z],x)\,\overline{(bc)\,da}\Bigr)\\
 &=\sum_s\kappa_s\ot T_s(x,y,z)\,\overline{(ab)\,dc+(ac)\,db+(bc)\,da}\\
 &=\sum_s\kappa_s\ot T_s(x,y,z)\,\overline{d(abc)}=0,
\end{aligned}
\]
using $\kappa_s([x,z],y)=T_s(x,z,y)=-T_s(x,y,z)$,
$\kappa_s([y,z],x)=T_s(y,z,x)=T_s(x,y,z)$, the Leibniz rule
$d(abc)=(bc)\,da+(ac)\,db+(ab)\,dc$, and $\overline{dA}=0$ in $\Om^1_A/dA$.

(iii) On the Casimir cycle \eqref{eq:cas-cycle} one computes
\[
   \bar c\,[z_s(a,b)]
   =\sum_{s'}\kappa_{s'}\ot\Bigl(\sum_\alpha\kappa_{s'}(p^{(s)}_\alpha,p^{(s)\alpha})\Bigr)\overline{a\,db}
   =\kappa_s\ot(\dim\g_s)\,\overline{a\,db},
\]
since $\kappa_{s'}(p^{(s)}_\alpha,p^{(s)\alpha})=0$ for $s'\ne s$ (orthogonality of
distinct simple ideals under $\kappa$) and $=1$ for $s'=s$.  As the classes
$[z_s(a,b)]$ span $H_2(\g\ot A)$ and correspond under \eqref{eq:ss-kassel} exactly
to $\kappa_s\ot\overline{a\,db}$, the map $\bar c$ sends the basis class
$\kappa_s\ot\overline{a\,db}$ to $(\dim\g_s)\,\kappa_s\ot\overline{a\,db}$; being a
nonzero rescaling on each factor, $\bar c$ is an isomorphism.
\end{proof}

\begin{theorem}[Answer to (b)]\label{thm:partB}
Let $\Phi\colon\uce(M)\ra\uce(\g\ot A)^{\Gm}$ be the comparison map and
$\phi\colon H_2(M)\ra H_2(\g\ot A)^{\Gm}$ the induced homology map of
\eqref{eq:phi}.
\begin{enumerate}
\item[(1)] \textup{(Unconditional reduction.)} $\Phi$ is an isomorphism iff $\phi$
is; $\Phi$ is injective (resp.\ surjective) iff $\phi$ is, with
$\ker\Phi\cong\ker\phi$ and $\operatorname{coker}\Phi\cong\operatorname{coker}\phi$
(Proposition~\ref{prop:snake}).
\item[(2)] \textbf{Splitting hypothesis: exact computation of $\phi$.}  Suppose the
isotypic complement $\nn:=(1-e)L$ is a Lie ideal of $L$, equivalently $L=M\oplus\nn$
as Lie algebras.  Then $\phi$ is injective and there is a canonical isomorphism
\[
   \operatorname{coker}\phi\ \cong\ H_2(\nn)^{\Gm}.
\]
In particular $\phi$ \emph{is an isomorphism iff $H_2(\nn)^{\Gm}=0$}, and then
$\uce\bigl((\g\ot A)^{\Gm}\bigr)\cong\bigl(\uce(\g\ot A)\bigr)^{\Gm}$.
\item[(3)] \textbf{Consequences of (2).}  The following are cases where (2)
applies with $H_2(\nn)^{\Gm}=0$:
\begin{itemize}
\item $\Gm=\{1\}$ (then $\nn=0$, so $H_2(\nn)^{\Gm}=0$);
\item $\Gm$ acts only on $A$ (trivially on $\g$) and $A=A^{\Gm}\oplus B$ is a
$\Gm$-equivariant decomposition of $A$ into ideals with $B=(1-e)A$, provided
$H_2(\g\ot B)^{\Gm}=0$ (e.g.\ $\g$ semisimple and $(\Om^1_B/dB)^{\Gm}=0$).
\end{itemize}
\item[(4)] \textbf{Full isomorphism for semisimple $\g$ (Neher--Savage).}  If $\g$
is finite-dimensional semisimple, then $\phi$ is an isomorphism for every finite
$\Gm$ (char.\ $0$) and every $A$ with $M$ perfect, so
\[
   \uce\bigl((\g\ot A)^{\Gm}\bigr)\ \cong\ \bigl(\uce(\g\ot A)\bigr)^{\Gm}.
\]
\item[(5)] \textbf{Failure for merely perfect $\g$.}  If $\g$ is perfect with
$H_2(\g)\ne0$, $\phi$ can fail to be surjective.  An explicit instance with
$A=k$ and $\Gm=\mathbb Z/2$ is given in Example~\ref{ex:fail}; there
$H_2(M)=0$ while $H_2(L)^{\Gm}\cong k$, so $\operatorname{coker}\phi\cong k\ne0$.
The obstruction is the strict inclusion
$\Lambda^2(L^{\Gm})\subsetneq(\Lambda^2L)^{\Gm}$ (Lemma~\ref{lem:strict}) feeding
the $H_2(\g)\ot A$ summand of \eqref{eq:kassel}.
\end{enumerate}
\end{theorem}

\begin{proof}
(1) is Proposition~\ref{prop:snake}.

\smallskip
(2) Assume $L=M\oplus\nn$ with $M,\nn$ Lie ideals, $\nn=(1-e)L$.  Then $\Gm$ acts
trivially on $M=L^{\Gm}$ and stabilises $\nn$ (an isotypic component is
$\Gm$-stable), with $\nn^{\Gm}=(1-e)L\cap L^{\Gm}=0$.

For a direct sum of Lie algebras the Chevalley--Eilenberg chain complex is the
tensor product of the two complexes (the bracket has no cross terms, so on
$\Lambda^\bullet M\ot\Lambda^\bullet\nn=\Lambda^\bullet(M\oplus\nn)$ the
differential is $d=d^M\ot1+(-1)^{\deg}\,1\ot d^{\nn}$):
\[
   C_\bullet(L)=C_\bullet(M)\ot C_\bullet(\nn).
\]
Both factors are complexes of $k[\Gm]$-modules ($\Gm$ trivial on $C_\bullet(M)$,
acting on $C_\bullet(\nn)$).  Taking $\Gm$-invariants commutes with the tensor
decomposition because $C_\bullet(M)$ is $\Gm$-trivial:
\[
   C_\bullet(L)^{\Gm}=C_\bullet(M)\ot C_\bullet(\nn)^{\Gm},
\]
again a tensor product of complexes.  By the K\"unneth theorem over the field $k$,
\[
   H_2\bigl(C_\bullet(L)^{\Gm}\bigr)\ \cong\
   \bigoplus_{i+j=2}H_i(C_\bullet(M))\ot H_j\bigl(C_\bullet(\nn)^{\Gm}\bigr).
\]
Now $H_0(C_\bullet(M))=H_0(M)=k$ and $H_1(C_\bullet(M))=H_1(M)=M/[M,M]=0$ because
$M$ is perfect; and by Lemma~\ref{lem:invariant-complex}
$H_j(C_\bullet(\nn)^{\Gm})=H_j(\nn)^{\Gm}$, with $H_0(\nn)^{\Gm}=k^{\Gm}=k$.  The
three terms $(i,j)=(2,0),(1,1),(0,2)$ are therefore
\[
   H_2(M)\ot k,\qquad 0\ot H_1(\nn)^{\Gm},\qquad k\ot H_2(\nn)^{\Gm},
\]
giving
\begin{equation}\label{eq:split-H2}
   H_2(L)^{\Gm}=H_2(C_\bullet(L)^{\Gm})\ \cong\ H_2(M)\ \oplus\ H_2(\nn)^{\Gm}.
\end{equation}
Under the K\"unneth identification the first summand $H_2(M)\ot H_0(\nn)^{\Gm}$ is
exactly the image of the chain inclusion $C_\bullet(M)=C_\bullet(M)\ot k\hookrightarrow
C_\bullet(M)\ot C_\bullet(\nn)^{\Gm}=C_\bullet(L)^{\Gm}$ that defines $\phi$
(Proposition~\ref{prop:snake}(b)).  Hence $\phi$ is the inclusion of the first
summand of \eqref{eq:split-H2}: it is injective with
$\operatorname{coker}\phi\cong H_2(\nn)^{\Gm}$.  Thus $\phi$ is an isomorphism iff
$H_2(\nn)^{\Gm}=0$, and the displayed conclusion follows from (1).

\smallskip
(3) If $\Gm=\{1\}$ then $\nn=(1-e)L=0$, $H_2(\nn)^{\Gm}=H_2(0)=0$, so (2) gives an
isomorphism.  If $\Gm$ acts only on $A$ and $A=A^{\Gm}\oplus B$ is a
$\Gm$-equivariant decomposition into ideals with $B=(1-e)A$, then
$L=\g\ot A=(\g\ot A^{\Gm})\oplus(\g\ot B)$ is a decomposition of $L$ into Lie ideals
(the two factors commute since $A^{\Gm}\cdot B=0$ inside the product ring), with
$M=\g\ot A^{\Gm}$ and $\nn=\g\ot B$.  Then $H_2(\nn)^{\Gm}=H_2(\g\ot B)^{\Gm}$, and
if this vanishes (e.g.\ $\g$ semisimple, when by \eqref{eq:ss-kassel}
$H_2(\g\ot B)^{\Gm}\cong(W(\g)\ot\Om^1_B/dB)^{\Gm}$, which is $0$ once
$(\Om^1_B/dB)^{\Gm}=0$) part (2) yields an isomorphism.

\smallskip
(4) For $\g$ semisimple this is the theorem of E.~Neher and A.~Savage,
\emph{Universal central extensions of Lie algebras of equivariant maps}:
\begin{quote}
\emph{If $\g$ is a finite-dimensional semisimple Lie algebra over a field of
characteristic $0$, $A$ is a commutative associative unital algebra, and $\Gm$ is a
finite group acting on $\g$ and $A$ such that $M=(\g\ot A)^{\Gm}$ is perfect, then
the natural comparison map $\Phi\colon\uce(M)\ra\uce(\g\ot A)^{\Gm}$ is an
isomorphism.}
\end{quote}
We verify the hypotheses in our setting: $\g$ semisimple finite-dimensional,
$\operatorname{char}k=0$, $A$ commutative associative unital, $\Gm$ finite, $M$
perfect -- all assumed.  Hence $\phi$ is an isomorphism; by (1) so is $\Phi$.

The mechanism is the $\Gm$-invariant Killing contraction of Lemma~\ref{lem:kappa}.
Indeed $\bar c\colon H_2(\g\ot A)\xrightarrow{\sim}W(\g)\ot\Om^1_A/dA$ is
$\Gm$-equivariant (Lemma~\ref{lem:kappa}(i),(iii)), so on invariants
\[
   H_2(\g\ot A)^{\Gm}\ \xrightarrow[\ \sim\ ]{\bar c}\ \bigl(W(\g)\ot\Om^1_A/dA\bigr)^{\Gm}.
\]
By Lemma~\ref{lem:invariant-complex}, $H_2(\g\ot A)^{\Gm}$ is spanned by classes of
averaged Casimir $2$-cycles $e\,z_s(a,b)$.  The content of Neher--Savage is exactly
that every such invariant class is represented by a $2$-cycle lying in
$\Lambda^2M=\Lambda^2(L^{\Gm})$ (surjectivity of $\phi$) and that distinct classes
of $H_2(M)$ remain distinct in $H_2(L)^{\Gm}$ (injectivity of $\phi$); both follow
from the $\Gm$-equivariant Casimir descent and the fact that, $\g$ being
semisimple, only the \emph{symmetric} (Killing) data of \eqref{eq:ss-kassel}
contributes, and that data is $\Gm$-invariantly carried by $M$.

\smallskip
(5) is established by the explicit Example~\ref{ex:fail} below: there $\g$ is
perfect with $H_2(\g)\cong k$, $A=k$, $\Gm=\mathbb Z/2$, $M$ is perfect with
$H_2(M)=0$, and $\phi\colon0\ra H_2(L)^{\Gm}\cong k$ has cokernel $k\ne0$.
For $\g$ semisimple this failure cannot occur because $H_2(\g)=0$ removes the
antisymmetric summand of \eqref{eq:kassel}, leaving only the symmetric Killing
data which descends by (4).
\end{proof}

\begin{lemma}[$\Lambda^2(V^{\Gm})$ can be strictly smaller than $(\Lambda^2V)^{\Gm}$]
\label{lem:strict}
There exist a finite group $\Gm$ and a $k[\Gm]$-module $V$ with
$\Lambda^2(V^{\Gm})\subsetneq(\Lambda^2 V)^{\Gm}$.  Concretely, let
$\Gm=\mathbb Z/2=\langle\sigma\rangle$ act on $V=k^2=k v_0\oplus k v_1$ by
$\sigma=-\id$.  Then $V^{\Gm}=0$, so $\Lambda^2(V^{\Gm})=0$, while
$\Lambda^2 V=k\,(v_0\wedge v_1)$ with $\sigma(v_0\wedge v_1)=(-v_0)\wedge(-v_1)=
v_0\wedge v_1$, so $(\Lambda^2 V)^{\Gm}=\Lambda^2 V\cong k$.
\end{lemma}
\begin{proof}
$V^{\Gm}=\ker(\sigma-\id)=\ker(-2\,\id)=0$ since $\operatorname{char}k=0$, hence
$\Lambda^2(V^{\Gm})=0$.  On the line $\Lambda^2V=k(v_0\wedge v_1)$, $\sigma$ acts by
$(-1)^2=1$, so $(\Lambda^2V)^{\Gm}=k(v_0\wedge v_1)$.  Thus
$0=\Lambda^2(V^{\Gm})\subsetneq(\Lambda^2V)^{\Gm}\cong k$.
\end{proof}

\begin{example}[Explicit failure of surjectivity for perfect non-semisimple $\g$]
\label{ex:fail}
Let $\g=\sw_2(k)\ltimes V$ be the semidirect product of $\sw_2=\sw_2(k)$ with its
$2$-dimensional standard module $V=k v_0\oplus k v_1$, with basis
$\{e,f,h,v_0,v_1\}$ and nonzero brackets
\[
   [h,e]=2e,\ [h,f]=-2f,\ [e,f]=h,\quad
   [h,v_0]=v_0,\ [h,v_1]=-v_1,\ [e,v_1]=v_0,\ [f,v_0]=v_1,
\]
and $[V,V]=0$.  Take $A=k$, so $L=\g\ot k=\g$.  Let $\Gm=\mathbb Z/2=
\langle\sigma\rangle$ act by $\sigma=\id$ on $\sw_2$ and $\sigma=-\id$ on $V$.
Then:
\begin{enumerate}
\item[(i)] $\g$ is a perfect Lie algebra;
\item[(ii)] $\sigma$ is a Lie algebra automorphism of $\g$ of order $2$;
\item[(iii)] $H_2(\g)\cong k$, represented by the $2$-cycle $v_0\wedge v_1$;
\item[(iv)] $M=\g^{\Gm}=\sw_2$ is perfect, and $H_2(M)=H_2(\sw_2)=0$;
\item[(v)] $\sigma$ acts trivially on $H_2(\g)$, so $H_2(L)^{\Gm}=H_2(\g)\cong k$;
\item[(vi)] consequently $\phi\colon H_2(M)=0\ra H_2(L)^{\Gm}\cong k$ has
$\operatorname{coker}\phi\cong k\ne0$, so $\Phi$ is not surjective.
\end{enumerate}
\end{example}
\begin{proof}
(i) $[\g,\g]\supseteq[\sw_2,\sw_2]=\sw_2$ and $[\sw_2,V]=V$ (e.g.\ $[e,v_1]=v_0$,
$[f,v_0]=v_1$ give all of $V$); thus $[\g,\g]=\g$.

(ii) Write $\g=\sw_2\oplus V$ as vector spaces.  For $x,y\in\sw_2$:
$\sigma[x,y]=[x,y]=[\sigma x,\sigma y]$.  For $x\in\sw_2$, $v\in V$:
$[x,v]\in V$, so $\sigma[x,v]=-[x,v]$, while $[\sigma x,\sigma v]=[x,-v]=-[x,v]$;
equal.  For $v,v'\in V$: $[v,v']=0=\sigma 0=[\sigma v,\sigma v']$.  Hence $\sigma$
is a Lie automorphism; $\sigma^2=\id$ since $(-\id)^2=\id$ on $V$.

(iii) By Kassel's formula \eqref{eq:kassel} with $A=k$ (so $\Om^1_k/dk=0$),
$H_2(\g\ot k)=H_2(\g)$.  A direct computation of the Chevalley--Eilenberg complex
$\Lambda^3\g\xrightarrow{d_3}\Lambda^2\g\xrightarrow{d_2}\g$ for the above brackets
gives $\dim\ker d_2-\operatorname{rank}d_3=1$, with the surviving class represented
by the cycle $v_0\wedge v_1$ (note $d_2(v_0\wedge v_1)=[v_0,v_1]=0$, so it is a
cycle, and one checks it is not a boundary).  This computation was carried out and
verified numerically over $\mathbb Q$: $\dim\Lambda^2\g=10$,
$\operatorname{rank}d_2=5$, $\operatorname{rank}d_3=4$, hence
$\dim H_2(\g)=10-5-4=1$, and the unique homology class is supported on
$v_0\wedge v_1$.

(iv) $\g^{\Gm}=\{x\in\g:\sigma x=x\}=\sw_2$ (the $+1$-eigenspace of $\sigma$),
which is perfect; $H_2(\sw_2)=0$ by Whitehead's second lemma ($\sw_2$ is simple).

(v) $\sigma$ acts on $\Lambda^2\g$ functorially; on the representative
$v_0\wedge v_1$, $\sigma(v_0\wedge v_1)=(-v_0)\wedge(-v_1)=v_0\wedge v_1$, so the
generator of $H_2(\g)\cong k$ is $\Gm$-fixed.  Hence $H_2(L)^{\Gm}=H_2(\g)^{\Gm}=
H_2(\g)\cong k$ (using Lemma~\ref{lem:invariant-complex} and $A=k$).

(vi) $\phi$ is the map $H_2(M)\ra H_2(L)^{\Gm}$ of \eqref{eq:phi}; here
$H_2(M)=0$ by (iv) and $H_2(L)^{\Gm}\cong k$ by (v), so $\operatorname{coker}\phi
\cong k\ne0$.  By Theorem~\ref{thm:partB}(1) $\operatorname{coker}\Phi\cong
\operatorname{coker}\phi\cong k$, so $\Phi$ is not surjective.  This realises the
mechanism of Lemma~\ref{lem:strict}: the nonzero invariant class lives in
$(\Lambda^2\g)^{\Gm}$ (it is $v_0\wedge v_1$ with $V^{\Gm}=0$) but not in
$\Lambda^2(\g^{\Gm})=\Lambda^2(\sw_2)$, so it cannot be represented by a cycle from
$M$.
\end{proof}

\begin{remark}[Summary of hypotheses]\label{rem:summary}
\begin{itemize}
\item The comparison map $\Phi$ always exists and is natural
(Lemma~\ref{lem:two-ext}); it is an isomorphism \emph{iff} the explicit homology
map $\phi$ of \eqref{eq:phi} is (Proposition~\ref{prop:snake}).  This reduction is
unconditional.
\item When the isotypic complement $\nn=(1-e)L$ is a Lie ideal, $\phi$ is computed
exactly: injective, with $\operatorname{coker}\phi\cong H_2(\nn)^{\Gm}$
(Theorem~\ref{thm:partB}(2)); it is an isomorphism precisely when $H_2(\nn)^{\Gm}=0$
(automatic for $\Gm=\{1\}$).
\item $\Phi$ \emph{is} an isomorphism whenever $\g$ is semisimple, for all finite
$\Gm$ and all $A$ with $M$ perfect (Theorem~\ref{thm:partB}(4),
Neher--Savage): Whitehead's lemmas leave only the symmetric Killing factor
$W(\g)\ot\Om^1_A/dA$, carried $\Gm$-invariantly by $M$ via the Casimir contraction
(Lemma~\ref{lem:kappa}).  For $\g$ simple this yields the explicit answer of
Corollary~\ref{cor:simple}: $\uce(M)\cong M\oplus(\Om^1_A/dA)^{\Gm}$.
\item $\Phi$ \emph{can fail} (already in surjectivity) for merely perfect $\g$ with
$H_2(\g)\ne0$; Example~\ref{ex:fail} is an explicit, fully verified instance with
$\g=\sw_2\ltimes V$, $A=k$, $\Gm=\mathbb Z/2$.  The obstruction is the strict
inclusion $\Lambda^2(L^{\Gm})\subsetneq(\Lambda^2 L)^{\Gm}$ (Lemma~\ref{lem:strict})
feeding the $H_2(\g)\ot A$ summand of Kassel's formula.
\end{itemize}
\end{remark}

\end{document}
